% ============================================================
\section{Introduction}
\label{sec:intro}
% ============================================================

Large language models have made open-ended educational dialogue practical, yet most deployed LLM tutors remain session-bounded assistants. They adapt to the immediate prompt rather than to a persistent learner model, and their functions are often implemented as isolated modules~\citep{chu-etal-2025-llm,letourneau2025systematicreviewaidriven}.
While a system may answer questions or generate exercises, these components share little state with one another. The result is not merely weaker personalization but a fragmented learning experience where pedagogical progress in one interaction rarely informs the next, as illustrated in Figure~\ref{fig:compare}.

Adopting a systems view, we argue that personalized tutoring requires a unified architecture rather than isolated prompts. This architecture must synchronize learner memory, task decomposition, and tool usage across multimodal sessions. Crucially, it should also provide a scalable framework for evolving from reactive to proactive support, obviating the need to rebuild the core runtime when introducing new interfaces or behaviors.

\begin{figure}[htbp]
    \centering
    \includegraphics[width=\columnwidth]{figures/compare-0411.png}
    \caption{Comparison between prior LLM products with \textsc{DeepTutor}. Conventional LLM tutors rely on generic priors, producing answers or exercises that may mismatch the learner's syllabus and proficiency, while \textsc{DeepTutor} grounds both problem solving and question generation in the learner's knowledge base and diagnosed weaknesses.}
    \label{fig:compare}
\end{figure}

Agent-Native Design. The agentic paradigm~\citep{yao2023reactsynergizingreasoningacting,wang2024survey} is uniquely suited to this challenge for three reasons.
First, agents can maintain persistent state across extended horizons, which is essential for cultivating a dynamic learner model rather than treating interactions as isolated events.
For instance, a tutoring agent can diagnose recurring misconceptions, such as a student's persistent confusion between the chain and product rules, and adaptively refine its instructional strategy when relevant topics re-emerge in future sessions.
Second, agents can orchestrate tools and sub-agents into structured workflows, enabling complex tasks like problem-solving, targeted exercise generation, and guided learning to be implemented as systematic pipelines rather than brittle prompt chains.
Third, the underlying agent infrastructure provides a scalable foundation for autonomous operation, allowing proactive behaviors to extend the system’s utility without re-engineering its core logic.
Consequently, we architect \textsc{DeepTutor} along two complementary dimensions: \emph{personalized agentic tutoring}, which establishes the core pipelines and the unifying personalization engine, and \emph{proactive autonomous companionship}, which leverages those pipelines for long-term, cross-platform engagement.
\begin{itemize}[leftmargin=*]
\item \textbf{Personalized Agentic Tutoring.}\quad
A hybrid personalization engine synthesizes static knowledge grounding with a dynamic \emph{trace forest}, distilling interaction history into a continuously refined learner profile.
We establish a closed-loop pedagogy that interlocks citation-grounded problem solving with difficulty-calibrated question generation. This shared substrate further powers collaborative writing, multi-agent deep research, and interactive guided learning, ensuring that pedagogical progress transcends individual modalities.

\item \textbf{Proactive Autonomous Companionship.}\quad
Through \textsc{TutorBot}, \textsc{DeepTutor} operationalizes its tutoring core into autonomous agents equipped with extensible skills, multi-bot coordination, and context-preserving multi-channel gateways.
This architecture transforms the system from a reactive interface into a proactive companion capable of autonomously initiating review sessions, synthesizing daily practice, and remediating diagnosed knowledge gaps across platforms.
\end{itemize}

This report presents \textsc{DeepTutor}, an open-source framework that actualizes this vision. The system is built upon a unified runtime where personalization, context propagation, and event-based streaming are natively shared across all workflows, spanning both reactive tutoring and proactive deployment. We offer this two-dimensional design as a generalizable blueprint for architecting personalized tutoring systems that seamlessly evolve into context-aware, proactive agents.

To evaluate the system’s core, we construct \textsc{TutorBench}, a student-centric benchmark comprising university-level materials across five diverse disciplines. Each entry synthesizes a source-grounded learner profile with diagnosed knowledge gaps and a corresponding interactive tutoring task. A first-person student simulator conducts multi-turn dialogues to rigorously assess adaptive pedagogical behaviors from the learner's perspective. 
In this report, we focus our quantitative evaluation on the foundational tutoring pipeline while characterizing cognitive extensions and \textsc{TutorBot} as system components whose long-term efficacy warrants extended longitudinal studies.

Our main contributions are as follows:
\begin{itemize}[nosep,leftmargin=*]
\item We propose \textsc{DeepTutor}, a novel intelligent tutoring framework that decouples a validated pedagogical core (\emph{personalized agentic tutoring}) from its deployment-oriented extensions (\emph{proactive autonomous companionship}), both anchored by a unified personalization engine.
\item We introduce a dynamic multi-resolution memory architecture that distills multi-turn interaction history into a continuously refined learner profile, enabling a tightly coupled bidirectional loop between problem solving and personalized question generation.
\item We design \textsc{TutorBot}, a proactive multi-agent layer that operationalizes tutoring capabilities through extensible skills, multi-agent coordination, and context-unified access across multiple platforms.
\item We construct \textsc{TutorBench}, a student-centric benchmark featuring a first-person interactive protocol to assess adaptive pedagogy. Empirical results demonstrate that \textsc{DeepTutor} improves personalized tutoring quality by 10.8\% over the strongest baseline, while simultaneously achieving an average 28.6\% gain in general reasoning across five backbone models on standard benchmarks.
\end{itemize}